\item Considere una muestra de $1.00\text{ kg}$ de uranio natural compuesto por $^{238}\text{U}$ ($99.275\%$ por masa), $^{235}\text{U}$ ($0.720\%$ por masa) y $^{234}\text{U}$ ($0.005\%$ por masa, $T_{1/2} = 2.44 \times 10^5\text{ años}$).
\begin{enumerate}
    \item Encuentre la actividad radiactiva en Curies debida a cada uno de los tres isótopos por separado.
    \item ¿Qué fracción de la actividad total de la muestra se debe a cada isótopo?
    \item A partir de sus resultados, discuta si la radiación de esta muestra representa un peligro de exposición agudo o crónico.
\end{enumerate}